Our workload consists of Astropy instances from the SWE-bench Verified test split. The agent uses Qwen Code 0.19.4 to inspect repositories, edit files, execute tools, and produce patches. A shared two-task comparison covers three tree runs and one SGLang EAGLE run that each produced nonempty patches on both tasks. A separate ten-task segment, Cqc10, recorded on 24 August 2026, supplies the latest complete single-configuration deployment case. The task/evaluation worker is x86-64; model serving runs on the GB10. These are full agent attempts. Table~\ref{tab:swe-tasks} lists every task and saved evaluator outcome in Cqc10.

\paragraph{Ten-task deployment}
The patched vLLM engine serves the RadixArk Qwen3.8-27B NVFP4 checkpoint with ModelOpt mixed quantization and bfloat16 execution dtype. The tree verifier uses the geometry and attention configuration specified in Section~\ref{sec:protocol}, with full-vocabulary drafting and reused draft logits. Prefix caching and full/piecewise CUDA graphs are enabled. The engine permits one active sequence, with maximum context 131,072 tokens and a 4,096-token scheduler budget. Source commit \texttt{e7af6b595a8b} and the loaded extension hash are retained in the artifact. The implementation uses the graph/cache-enabled stochastic route described above.

The proxy supplies temperature 0.6, top-$p$ 0.95, top-$k$ 20, minimum-$p$ 0, and presence penalty 1.0; the engine seed is zero. The per-response output cap is 24,000 tokens, with a separate 20,000-token compaction cap and a 9,000-second agent timeout. Proxy auto-continuation is disabled. The agent runs in the task's repository image on an internal network that permits the model proxy and blocks external egress; per-task records retain the verified network-policy fingerprint. Agent duration comes from the task runner's elapsed-time field: it includes the agent loop and tool work, and excludes server startup and subsequent evaluation. We do not infer a completed-task speedup from model-only counters.

\begin{table}[t]
\caption{All ten tasks in the completed Cqc10 Astropy case. IDs have prefix \texttt{astropy\_\_astropy-}. Times are agent elapsed minutes, rounded from saved runner metadata; evaluation is subsequent. This selected development segment has one attempt per listed task and no matched native arm.}
\label{tab:swe-tasks}
\centering\small
\begin{tabular}{@{}lrl@{}}
\toprule
Task suffix & Agent time (min) & Saved evaluation \\
\midrule
13977 & 24.59 & Failed tests \\
14096 & 22.99 & Resolved \\
14182 & 7.23 & Failed tests \\
14309 & 1.67 & Resolved \\
14365 & 6.78 & Failed tests \\
14369 & 27.70 & Resolved \\
14508 & 41.24 & Resolved \\
14539 & 18.57 & Resolved \\
14598 & 25.79 & Failed tests \\
14995 & 4.59 & Resolved \\
\bottomrule
\end{tabular}
\end{table}

\paragraph{Workload observations}
All ten attempts have evaluator reports: six resolved their tasks and four failed tests. Recorded agent durations sum to 181.15 minutes, with a range of 1.67--41.24 minutes. This is an observed task ledger, not a representative SWE-bench Verified success estimate. The sum is accumulated agent time, not elapsed campaign time including startup and evaluation.

Every task used tools and produced a nonempty patch. Analysis of the recorded traces detects no degeneration in these ten runs and no malformed tool-call arguments. The degeneration heuristic jointly requires at most 200 visible characters, zero tool calls, and at least 2,000 thinking characters; none of the ten traces meets that conjunction. This screen can miss loops that continue narrating or using tools, so its zero detections establish neither universal degeneration freedom nor preserved task quality. The earlier campaign's detected degeneration remains an adverse observation in the lineage audit.

\paragraph{Pooled decode comparison}
For completed request $r$, let $n_r$ be its output-token count, $\ell_r$ its server-recorded end-to-end latency, and $f_r$ its time to first output. With $R$ completed requests and $N=\sum_r n_r$, we use
\begin{equation}
 D_{\mathrm{pool}} =
 \frac{\sum_{r=1}^{R}(n_r-1)}{\sum_{r=1}^{R}(\ell_r-f_r)}
 = \frac{N-R}{\sum_r\ell_r-\sum_r f_r}.
 \label{eq:pooled-decode}
\end{equation}
Counts and durations are pooled before division. The numerator uses the conventional $n_r-1$ output intervals; outputs arriving in one batch can have zero arrival gaps. The denominator excludes first-output latency and between-request agent/tool work. It is accumulated request time, which counts overlapping requests separately, not isolated GPU time or elapsed campaign time.

For Cqc10, $N=248{,}342$, $R=265$, $\sum_r\ell_r=10{,}667.6211$ seconds, and $\sum_r f_r=989.8239$ seconds yield 25.63 tokens/s. The 265 requests reconcile with 256 normal requests and nine successful internal compaction requests; all recorded completion reasons are stop. The ten pre/post metric brackets are contiguous. Generation-token and completed-request counters agree with their latency-histogram populations, with no counter reset or outstanding request at the selected boundaries.

The eligible SGLang EAGLE comparator, recorded on 24 August, serves the as-shipped Qwen3.8-27B NVFP4 checkpoint with FlashInfer, using speculative steps/top-$k$/draft-token settings of 3/1/4. It uses the same Qwen Code version, agent harness, requested sampling settings, and network policy, with a 24,000-token response cap, a 262,144-token context limit, and an 8,192-token prefill chunk. Its two completed Astropy tasks are 12907 (resolved; 7.10 agent minutes) and 13033 (failed tests; 26.96 agent minutes). Their counters give $N=47{,}854$, $R=45$, $\sum_r\ell_r=2{,}013.6344$ seconds, and $\sum_r f_r=235.9488$ seconds, yielding 26.89 tokens/s. Both brackets contain streaming requests only. All completed engine requests are included, including internal agent traffic.

\paragraph{Performance eligibility}
After inspecting the archived runs, we restrict the shared-task rate comparison to run-pairs in which both tasks completed and produced nonempty patches. We apply this retrospective rule to both systems and retain patches that fail tests. All three tree pairs and one of the two recovered SGLang pairs meet the rule. The other SGLang pair is excluded: task 13033 produced an empty patch after a final thinking-only response at the configured output cap, so the test harness was not invoked. Its unfiltered pair rate is 30.70 tokens/s; Table~\ref{tab:all-attempts} retains that rate and failed outcome as a diagnostic. The earlier pair used a 32,768-token response cap, so this failure cannot be attributed to the later comparator's smaller 24,000-token cap. We exclude the entire pair from this performance comparison, rather than retaining only its successful task. The rates therefore characterize patch-producing runs, not performance over all attempted tasks.

\paragraph{Exploratory shared-task observations}
We select task IDs 12907 and 13033 because both have completed SGLang records, then report all three recovered split-K tree runs on that pair. Sr12 and Cqc16 ran on 19 August; Cqc15 ran on 23 August. All three use the same tree geometry, full-vocabulary drafting, and grouped split-K attention configuration as Cqc10. Runtime checks confirm the specialized attention kernel in all sixteen full-attention layers, with no fallback. They share the loaded attention extension but differ in Python patch revisions. Each uses a 32,768-token response cap, compared with the eligible SGLang run's 24,000. All 41 Sr12 responses and all 45 eligible SGLang responses are at most 20,000 tokens, below both configured caps; no response uses the differing output allowance. All Sr12 requests terminate with stop.

The August 19 revision \texttt{bdca0bd} records $N=25{,}027$, $R=41$, $\sum_r\ell_r=989.2612$ seconds, and $\sum_r f_r=130.4422$ seconds give 29.09 tokens/s under Eq.~\ref{eq:pooled-decode}. Other revisions give 28.28 and 27.27; Table~\ref{tab:pooled-decode} retains each observation. Selecting the maximum of these revisions does not estimate fixed-build repeatability or generalization. All selected task brackets pass the same population and idle-boundary checks. Every row on the shared pair has one resolved task and one failed task. All six selected tree attempts used tools and produced nonempty patches, with no detected degeneration or malformed tool-call arguments under the recorded checks. Every selected tree request finished with stop, without a recorded response-cap hit.

\begin{table}[t]
\caption{Pooled decode rates for tree verification and SGLang EAGLE with Qwen3.8-27B NVFP4 on GB10, using Eq.~\ref{eq:pooled-decode}. The first four rows use the same two SWE-bench Verified Astropy tasks, 12907 and 13033; each included run produces nonempty patches on both tasks; decoder settings are specified in the text. The August 24 tree deployment covers a separate ten-task subset. Revisions are exploratory observations, not statistical replicates.}
\label{tab:pooled-decode}
\centering\small
\begin{tabular}{@{}lrrr@{}}
\toprule
Deployment & Tasks & Requests & Tokens/s \\
\midrule
SGLang EAGLE & 2 & 45 & 26.89 \\
Tree bdca0bd (Aug 19) & 2 & 41 & 29.09 \\
Tree 78a29d3 (Aug 19) & 2 & 35 & 28.28 \\
Tree a450f6c (Aug 23) & 2 & 53 & 27.27 \\
\midrule
Tree e7af6b5 (Aug 24) & 10 & 265 & 25.63 \\
\bottomrule
\end{tabular}
\end{table}

The shared-pair result is a retrospective subset of the recorded workloads. Sr12 completed four tasks in total, two resolved and two failed; pooling all four gives 26.21 tokens/s. Cqc16's subsequent task 13236 exhibited degeneration, and Cqc15's subsequent task 13398 lacks a completed evaluation. Those adverse records remain in the artifact and outside this explicitly two-task comparison. The shared-task comparison evaluates the complete decoding systems, including their proposal topology, draft budget, and serving optimizations. These method-specific choices are part of the comparison. The records use the same requested sampling settings and agent harness. The performance table retains every eligible run-pair: three tree observations and one SGLang observation. This retrospective subset does not establish performance over all attempts or each system's optimum over draft lengths and candidate budgets. Repeated runs with a predeclared tuning protocol would measure that performance envelope; a component ablation would separately isolate the contribution of tree verification. The artifact retains the original metric brackets, source-bound configuration and behavior records, and reproducible pooled reductions.

\paragraph{All-attempt outcome accounting}
The recovered ledger has 22 unique attempts with complete task metric brackets: four in the August 19 \texttt{bdca0bd} run, two in each other shared-pair tree run, ten in the August 24 deployment, and two in each SGLang pair. Table~\ref{tab:all-attempts} retains their outcomes and pooled rates, including the empty-patch SGLang pair. The two-task \texttt{bdca0bd} row is contained in its four-task row; the rows must not be summed as independent cohorts. Additional partial and unbracketed records stay in a separate provenance ledger, without invented rates or completed evaluations.

\begin{table*}[t]
\caption{All recovered complete-bracket development cohorts. Rates are descriptive diagnostics under Eq.~\ref{eq:pooled-decode}, including failed attempts. The earlier SGLang pair ($^*$) contains an empty patch and is excluded only from the patch-producing comparison in Table~\ref{tab:pooled-decode}. The first tree pair overlaps the final four-task row; different task populations are not ranked against one another.}
\label{tab:all-attempts}
\centering\small
% Generated by reduce_p0.py. Exploratory cohorts; Sr12 subset overlaps full-four cohort.
\begin{tabular}{lrrrr}
\toprule
Build/cohort & Tasks & Resolved & Agent min & Pooled tok/s \\
\midrule
Aug. 19 tree (bdca0bd) & 2 & 1 & 17.31 & 29.09 \\
Aug. 19 tree (78a29d3) & 2 & 1 & 18.06 & 28.28 \\
Aug. 23 tree (a450f6c) & 2 & 1 & 48.68 & 27.27 \\
Earlier SGLang$^{*}$ & 2 & 1 & 26.70 & 30.70 \\
Later SGLang & 2 & 1 & 34.07 & 26.89 \\
Aug. 24 tree (10 tasks) & 10 & 6 & 181.15 & 25.63 \\
Aug. 19 tree (4 tasks) & 4 & 2 & 73.67 & 26.21 \\
\bottomrule
\end{tabular}
% * Earlier SGLang remains visible as an all-attempt diagnostic: one empty patch; not an eligible patch-producing pair.

\end{table*}

\begin{table*}[t]
\caption{Shared-task agent times and outcomes, including the excluded earlier SGLang pair ($^*$). Task suffixes identify Astropy instances 12907 and 13033. Agent time includes model service and tools, but excludes subsequent evaluation. Different token trajectories are outcomes of the systems and do not isolate kernel cost.}
\label{tab:paired-agent-time}
\centering\footnotesize
% Generated shared-task rows; keeps failed/empty-patch attempts.
\begin{tabular}{llrrrl}
\toprule
Build & Task & Agent s & Output tokens & Patch bytes & Verdict \\
\midrule
Aug. 19 tree (bdca0bd) & 12907 & 407.188 & 7873 & 504 & resolved \\
Aug. 19 tree (bdca0bd) & 13033 & 631.567 & 17154 & 1017 & failed \\
Aug. 19 tree (78a29d3) & 12907 & 291.437 & 5603 & 504 & resolved \\
Aug. 19 tree (78a29d3) & 13033 & 791.870 & 20636 & 1092 & failed \\
Aug. 23 tree (a450f6c) & 12907 & 670.872 & 14138 & 504 & resolved \\
Aug. 23 tree (a450f6c) & 13033 & 2249.684 & 56286 & 1450 & failed \\
Earlier SGLang$^{*}$ & 12907 & 238.715 & 5895 & 504 & resolved \\
Earlier SGLang$^{*}$ & 13033 & 1363.453 & 41955 & 0 & failed \\
Later SGLang & 12907 & 426.131 & 8314 & 504 & resolved \\
Later SGLang & 13033 & 1617.769 & 39540 & 912 & failed \\
\bottomrule
\end{tabular}

\end{table*}

Table~\ref{tab:paired-agent-time} supplies the missing task-level timing context. The August 19 \texttt{bdca0bd} pair totals 1,038.755 agent seconds; later SGLang totals 2,043.900 seconds. Both resolve task 12907 and fail task 13033. The token totals differ substantially. These single recorded pairs cannot establish a general task-completion speedup, and the spread across tree revisions cannot estimate within-build uncertainty.


\paragraph{Selection and comparison limits}
Cqc10 contains the ten remaining tasks of an interrupted sixteen-task development campaign. Earlier segments changed source revisions and output caps and included degeneration, incomplete output, and budget-capped attempts. Source-bound audit records of those failures, retries, and missing evaluations remain in the artifact; the original records remain at their recorded source paths. We do not relabel the completed segment as the entire campaign, combine those segments into a single-configuration benchmark score, or drop unsuccessful attempts from a claimed benchmark denominator. The tasks were selected from one repository, and this deployment has one recorded boot without a matched ten-task native run. Separate targeted native probes change other serving settings and do not supply that comparator. Consequently, these records demonstrate an implemented verifier serving real coding-agent work, but establish no native-versus-tree task-speed or quality advantage. A current, matched task campaign is the remaining experiment for that claim.


\paragraph{Existing comparator experiments}
Earlier SGLang EAGLE trials and native MTP controls remain in the repository. Appendix~\ref{app:history} recovers all 48 task outcomes from a July FP8 tree/MTP-5/MTP-11 campaign, with ten, ten, and eight resolved tasks out of the same sixteen, respectively. That superseded design does not provide a native control for the NVFP4 deployment; later NVFP4 native probes cover targeted tasks with different serving settings. DSpark measurements use synthetic requests, and a later DFlash2 check establishes boot and response handling rather than agent-task performance. These records are distinguished from the pooled SWE comparison above; their superseded or non-agent rates are not additional performance results for the current verifier.
